% paper/sections/09b1_split_ppe.tex
%
% §9.3  分相 PPE 解法：密度比非依存な高精度 PPE 戦略
% ═══════════════════════════════════════════════════════════════════════════════
\subsection{分相 PPE 解法：密度比非依存な高精度 PPE 戦略}
\label{sec:split_ppe_framework}
% backward-compat labels
\label{sec:split_ppe_motivation}%
\label{sec:gfm_numerical_motivation}%
% ═══════════════════════════════════════════════════════════════════════════════

\noindent\textbf{本節の位置づけ：}
§~\ref{sec:ccd_poisson_matrix} で構築した CCD 離散化は
等密度ポアソン方程式に対して $\Ord{h^6}$ の設計精度を達成する．
しかし二相流の PPE は変密度演算子 $\bnabla\cdot(\rho^{-1}\bnabla p) = q$ であり，
界面を横切るステンシルが $\Ord{1}$ の打ち切り誤差を生じる．
本節では，この本質的制約を回避する\textbf{分相 PPE 解法}を本稿の主要な圧力閉包として定式化する．
一言でいえば，全領域を一本の変密度 Poisson 問題として扱わず，
液相・気相では定数密度 Poisson 問題を解き，界面では圧力ジャンプ条件で両相を接続する，
という分解である．
変密度圧力補正をそのまま変係数 Poisson 問題へ押し込まず，
圧力勾配項の分割により高速な定数係数 Poisson 解へ近づける発想は
Dodd--Ferrante~\cite{DoddFerrante2014} と問題意識を共有する．
ただし本稿の主眼は FFT 化ではなく，切断界面での圧力ジャンプ，HFE，DC 残差受理条件，
および毛管面補正を同じ面勾配契約へ載せることである．

\noindent\textbf{本節の閉包条件：}
分相 PPE は単一の離散化名ではなく，次の受理条件群を同時に満たす圧力閉包を指す．
\[
\begin{aligned}
  &\text{相内定数密度 Poisson}
  + j_{gl}=p_g-p_l=-\sigma\kappa_{lg}
  + G_\Gamma(p;j)=G(p)-B(j)\\
  &\qquad
  + \mathrm{HFE}
  + c_{\sigma,f}^{\mathrm{VW}} + H_{\mathrm{aug}}c_{\sigma,f}
  + \mathrm{DC}\{\eta_{\mathrm{DC}}\le\varepsilon_{\mathrm{PPE}}\}
  + \text{大域ゲージ・再投影分離}
\end{aligned}
\]
この閉包では，smoothed-Heaviside 一括 PPE は係数不連続の残るモデル，
CSF 体積力は Balanced--Force 原理を説明するモデル，
滑らかな補助圧力場によるジャンプ表現は代数的な対照モデルとして扱う．
いずれも，界面ジャンプを PPE と速度補正が共有する
面勾配条件へ直接入れる本節の閉包とは区別する．

% ─────────────────────────────────────────────────────────────────────────────
\subsubsection{smoothed-Heaviside 一括 PPE の本質的制約}
\label{sec:monolithic_ppe_limitation}
% ─────────────────────────────────────────────────────────────────────────────

smoothed-Heaviside 一括 PPE では，遷移層（幅 $\varepsilon = 1.5h$）内で
$1/\rho$ が $\rho_l/\rho_g$ 倍のジャンプを持つ．
界面を横切る CCD/有限差分ステンシルは $\Ord{1}$ の打ち切り誤差を生じ，
PPE のグリーン関数の非局所性により解全体に伝播する
（§~\ref{sec:ppe_condition_number}）．
この誤差は反復回数 $k$ の増大や緩和係数 $\omega$ では回復不能であり，
smoothed-Heaviside 一括 PPE の\textbf{本質的制約}である：
%
\begin{itemize}[noitemsep]
  \item 中程度密度比：高解像度側で格子収束がプラトー化
  \item 中-高密度比：頑健な DC 発散閾値は仮定せず，界面型密度場で設計精度は回復しない
  \item 水--空気級高密度比：smoothed-Heaviside 一括 PPE のみでの高精度閉包は未保証
\end{itemize}
%
したがって本章では，smoothed-Heaviside 一括 PPE の高密度比利用を理論保証外とする．
重要な点として，中程度の密度比でも界面近傍の精度は
CCD の設計精度 $\Ord{h^6}$ から\textbf{大幅に劣化}しており，
低密度比だから smoothed-Heaviside 一括 PPE で十分ということはない．

\noindent\textbf{条件数の定量的見積もり：}
smoothed-Heaviside 一括 PPE 演算子 $\bnabla\cdot(\rho^{-1}\bnabla\cdot)$ の条件数は
$\kappa_2 \approx (\rho_l/\rho_g)\cdot h^{-2}$ のオーダーで成長する．
水--空気級の高密度比（$\rho_l/\rho_g \sim 10^3$），$h = 1/64$ では
$\kappa_2 \sim 10^3\cdot 4096 \approx 4\times 10^{6}$ となり，
ゲージ射影後の非零固有空間でも欠陥補正反復の収束が遅い
（Neumann 零モードそのものは §~\ref{sec:ppe_phase_gauge_pin} で除去する）．
分相 PPE では各相内 $\kappa_2 \approx h^{-2}\sim 10^{3}$ に縮退する
（高密度比 PPE の解析と緩和戦略は~\cite{Guermond2006,GuermondQuartapelle1998}；
界面ジャンプを陽的に扱う Immersed Interface Method の系統は~\cite{LeVequeLi1994,Li2003IIMOverview}）．

% ─────────────────────────────────────────────────────────────────────────────
\subsubsection{分相 PPE 解法の定式化}
\label{sec:split_ppe_formulation}
% ─────────────────────────────────────────────────────────────────────────────

分相 PPE 解法の核心は「各相で定数密度のラプラシアンを解く」ことにある．
smoothed-Heaviside 一括 PPE が $\bnabla\cdot(\rho^{-1}\bnabla p) = q$ を
変動係数 $\rho(\psi)$ のまま全領域で解くのに対し，
分相 PPE 解法は各相 $k\in\{\ell,g\}$ で
%
\begin{equation}
  \nabla^2 p_k = \rho_k \, q \qquad \text{（$\rho_k$ は定数）}
  \label{eq:split_ppe}
\end{equation}
%
を解く（ここで $q = (1/\Delta t)\,\bnabla\cdot\bu^\ast$ は
予測速度由来の発散ソースであり，射影段階の右辺である）．
ただし，これは単に各相を切り離して解くという意味ではない．
界面 $\Gamma$ 上では圧力ジャンプと法線フラックス条件を同時に課して，
二つの相の楕円型問題を閉じる：
\begin{align}
  [p]_\Gamma
    &= J_p^\chi
     \quad \text{（未知量 $\chi$ に対応する圧力ジャンプ）},
     \label{eq:split_ppe_pressure_jump} \\
  \left[\frac{1}{\rho}\pd{p}{n}\right]_\Gamma
    &= J_n .
     \label{eq:split_ppe_flux_jump}
\end{align}
\noindent\textbf{全圧力と圧力増分の切り分け：}
以後の式では簡潔のため未知量を $p$ と書くが，
PPE に渡すジャンプは\textbf{解く未知量に対応するジャンプ}でなければならない．
すなわち
\begin{equation}
  J_p^\chi =
  \begin{cases}
    j_{gl}^{\,n+1}, & \chi = p^{n+1}\ \text{を直接解く場合},\\[2mm]
    \delta j_{gl}=j_{gl}^{\,n+1}-j_{gl}^{\,n},
      & \chi=\delta p=p^{n+1}-p^n\ \text{を解く場合}.
  \end{cases}
  \label{eq:split_ppe_unknown_jump_contract}
\end{equation}
この区別を失うと，Young--Laplace 圧力を二重計上するか，
逆に前時刻圧力の自由度に吸収して消してしまう．
本稿のジャンプ補正付き面勾配の記法で $j_{gl}$ と書く箇所は，
その射影ポテンシャルが保持すべき
$J_p^\chi$ を指す（全圧力を解く場合は $j_{gl}^{n+1}$，
圧力増分を解く場合は $\delta j_{gl}$）．
同様に法線速度の連続性から
$J_n=\Delta t^{-1}[\bu^*\cdot\hat{\bm n}_{lg}]$ を用いる．
予測速度が界面法線方向に連続となる保存的離散化では $J_n=0$ に退化する．

\noindent\textbf{$J_n$ の物理的正当化：}
連続 NS では界面で法線速度が連続（$[\bu\cdot\hat{\bm n}_{lg}]_\Gamma=0$；非相変化）．
予測速度の離散化が保存的ならば $\bu^*$ も界面で連続となり $J_n\equiv 0$．
非保存的離散化（例：面補間が界面で非対称）では
$[\bu^*\cdot\hat{\bm n}_{lg}]\sim \Ord{h^p}$ の残差が生じ，
$J_n$ は PPE でこの残差を吸収して速度補正後の $\bu^{n+1}$ を
連続に戻す役割を果たす（$\Ord{h^p}$ 一致時には $J_n=0$ とみなせる）．
HFE（§~\ref{sec:field_extension}）は，これらのジャンプ条件で補正した
片側データを界面越しに延長し，CCD ステンシルが不連続値に触れないようにするための
補助条件である．
$\rho_k$ が定数であるため，各相の離散化演算子は密度ジャンプの影響を一切受けない．

\noindent\textbf{欠陥補正法との関係：}
式~\eqref{eq:split_ppe} は等密度ポアソン方程式であるから，
§~\ref{sec:defect_correction_ppe} の欠陥補正法を各相内で適用できる．
等密度条件の下では DC の CCD 高次演算子 $L_H$ と有限差分低次演算子 $L_L$ の
スペクトル不整合が最小化されるため，
高次残差条件 $\eta_{\mathrm{DC}}\le\varepsilon_{\mathrm{PPE}}$ まで解けば
§~\ref{sec:ppe_dc_accuracy} の漸近精度が\textbf{密度比によらず}回復する．
反復回数そのものではなく，この残差条件へ到達することを物理閉包の条件とする．

\begin{tcolorbox}[resultbox, title={分相 PPE + DC 残差受理条件の相内精度と界面律速}]
\begin{enumerate}[noitemsep]
  \item 各相で定数密度の PPE を独立に解くため，
        界面条件~\eqref{eq:split_ppe_pressure_jump}--\eqref{eq:split_ppe_flux_jump}を満たす
        片側問題として閉じる
  \item DC は高次残差条件
        $\eta_{\mathrm{DC}}\le\varepsilon_{\mathrm{PPE}}$ で受理する．
        Dirichlet 製造解の成分検証では $\Ord{h^7}$ の超収束を確認し，
        Neumann 境界条件では $\Ord{h^5}$ が上界になる
  \item 各相内の離散誤差は密度比に直接依存しない．
        ただしこの主張は\textbf{相内部の Poisson 解}に対するものであり，
        界面切断面の閉包まで含めた全体精度を一律に $\Ord{h^6}$ と主張するものではない
  \item 速度補正の勾配 $\bnabla(\delta p)$ には CCD $\Ord{h^6}$（均一格子）
        または $\mathcal{G}^\text{adj}$（非一様格子 + 壁面境界条件；§~\ref{sec:fvm_ccd_corrector}）を適用する．
        この閉包では同じ面ジャンプ $B_f$ も補正側へ渡し，
        $G_\Gamma=G-B$ として Balanced--Force 整合性を保持する
  \item 界面近傍の最終律速は，FCCD 面ジェット
        （$u_f,u'_f$ は $\Ord{H^4}$，$u''_f$ は $\Ord{H^2}$；
        §~\ref{sec:face_jet_def}），HFE/GFM 延長，曲率評価の最小次数で決まる．
        したがって本節の高次性は
        \emph{相内 $\Ord{h^6}$ / DC 高次解 + 界面条件処理の別律速}
        と読むべきである
  \item 毛管補正量は式~\eqref{eq:capillary_surface_energy_riesz_cochain}
        の表面エネルギーの仮想仕事から構成し，
        式~\eqref{eq:capillary_hodge_gate} の毛管平衡判定で
        静的平衡成分と非平衡駆動成分を分ける
\end{enumerate}
\end{tcolorbox}

\noindent\textbf{HFE の役割：}
分相 PPE 解法において HFE（§~\ref{sec:field_extension}）は不可欠な場延長理論である．
ただし射影後の面速度を保存する圧力ジャンプ IPC では，
予測段階の履歴圧力項も節点スカラー $p^n$ の直接微分として扱わない．
前時刻で構成した物理圧力 $p^n$ と同じ界面ジャンプ $j^n$ から，
\begin{equation}
  \label{eq:affine_pressure_history_faces}
  \mathbf{a}^{\,n}_{p,f}
  =
  A_f\left(G_f p^n - B_f(j^n)\right)
\end{equation}
を履歴圧力面加速度として保存し，次の予測段階では
$\mathbf{u}^{*}_f$ に $-\Delta t\,\mathbf{a}^{\,n}_{p,f}$ として入れる．
すなわち HFE は相内の滑らかな延長を担うが，
切断面を横切る圧力履歴 / 補正段の離散条件は
式~\eqref{eq:affine_pressure_history_faces} と同じ
$G_\Gamma=G_f-B_f$ の面空間で閉じる．

% ─────────────────────────────────────────────────────────────────────────────
\subsubsection{圧力ジャンプ型閉包の範囲と純高次参照極限}
\label{sec:pure_fccd_split_ppe}
% ─────────────────────────────────────────────────────────────────────────────

本節の \textbf{圧力ジャンプ型閉包}は，
相内定数密度 PPE，ジャンプ補正付き面勾配，HFE，
表面エネルギー仮想仕事に基づく毛管平衡判定，DC 残差受理条件，
同一面フラックス整合，大域ゲージ・再投影分離を同時に満たす閉包である．
これは「全ての部品を純 FCCD だけで閉じる」ことと同義ではない．
DC の低次補正演算子，壁面閉包，面フラックス射影には
低次または有限体積法的な補助演算子が入り得るが，
PPE と速度補正が共有する面フラックス整合を破らない限り
同じ圧力ジャンプ型閉包に属する．

\noindent\textbf{純 FCCD 高次極限：}
純 FCCD 分相 PPE は有限体積法の制御体積保存を離散不変量として放棄し，
すべての空間演算子を FCCD / HFE / GFM の\textbf{単一高次言語}に統一する
高次極限である．質量保存は
$\Ord{\Delta x^4}$--$\Ord{\Delta x^6}$ の漸近的な高次性として扱われ，
制御体積恒等式としては保証されない．
したがって純 FCCD 高次極限は，本節の閉包の必要条件ではなく，
面演算子と誤差律速を説明する参照極限である．

\noindent\textbf{設計選択の根拠：}
有限体積法の制御体積保存を放棄する代償として，
(i) CLS の体積補正・質量閉包
（§\ref{sec:cls_stages} 段階 B と段階 F）が\textbf{体積質量保存}
（全領域 $\int\psi\,\mathrm{d}V$）を独立に担保する，
(ii) 速度場の離散発散ゼロ性は CCD 発散 $\bnabla\cdot\bu^*$ から構成される
PPE 右辺の射影で保証される，の二経路が保存性を補完する．
有限体積法では両者が同一の面フラックス恒等式から派生するのに対し，
本稿は「演算子の高次性」と「幾何学的発散ゼロ性」を独立系統で扱う．
数値検証では，質量保存と Balanced--Force 整合性を分離して評価する．

\noindent\textbf{数理構成の階層}：
\begin{enumerate}[leftmargin=2em,noitemsep]
  \item 界面幾何：Ridge--Eikonal（\S\ref{sec:ridge_eikonal}）
        で位相・計量を分離．
  \item 移流：FCCD + HFE（\S\ref{sec:fccd_advection}）で $\psi$ と
        運動量を共通面ジェット $\mathcal{J}_f$ 上で更新．
  \item 粘性：3 層構成（\S\ref{sec:viscous_3layer}）．
        $\mu\bnabla^2\bu$ は\textbf{界面越えでは使わない}．
  \item 圧力：分相 FCCD PPE + GFM ジャンプ行（本節）．
  \item 剛性処理：デフェクト補正
        （\S\ref{sec:defect_correction_ppe}）の高次外側 / 低次内側分解．
\end{enumerate}

% ─────────────────────────────────────────────────────────────────────────────
